Up to now, we have been considering negative results about \textbf{functions} as computational tasks. In particular, we defined a method to describe
which tasks are easy or difficult to solve, studying the border between the \textbf{feasible} and the \textbf{unfeasible}. \vspace{7pt}

Would it be possible to apply the same reasoning also for tasks coming from \textbf{machine learning}? Therefore, our purpose is to include inside what 
we have seen so far also functions related to machine learning tasks. We would like to understand which are the limits of machine learning, defining tasks 
achievable and not-achievable. \vspace{7pt}

Formally, this kind of extension is called \textbf{computational learning theory} and it gives us answers to the following questions:
\begin{itemize}
    \renewcommand{\labelitemi}{-}
    \item Which functions are easy to learn?
    \item Are there functions which are instead hard to learn?
    \item How much data is required to learn a given function?
    \item Can we learn functions with limited computational resources?
\end{itemize}